% Luc Destombe --- licence CC BY-NC-SA 4.0
% https://creativecommons.org/licenses/by-nc-sa/4.0/deed.fr
% Fichier autonome : compiler avec pdflatex, deux fois (signets, nombre de pages).
\documentclass[a4paper,10pt]{article}
\makeatletter
\newif\iflycee@niveaux
\lycee@niveauxtrue
\newif\iflycee@palatino
\lycee@palatinotrue

\RequirePackage[utf8]{inputenc}
\RequirePackage[T1]{fontenc}
\RequirePackage[french]{babel}
\RequirePackage{etoolbox}

\newcommand{\ShorthandsOff}{\@ifundefined{shorthandoff}{}{\shorthandoff{;:!?}}}
\newcommand{\ShorthandsOn}{\@ifundefined{shorthandon}{}{\shorthandon{;:!?}}}
\ShorthandsOff
\AtBeginDocument{\ShorthandsOn}
\AtBeginEnvironment{tikzpicture}{\ShorthandsOff}

\RequirePackage{geometry}
\RequirePackage{fancyhdr}
\RequirePackage{lastpage}
\RequirePackage{multicol}
\RequirePackage{array}
\RequirePackage{enumitem}
\RequirePackage{xcolor}
\RequirePackage{needspace}

\RequirePackage{amsmath,amssymb}
\RequirePackage{mathpazo}
\DeclareMathSizes{10.95}{10.95}{8}{5.5}
\begingroup\catcode`\_=\active \catcode`\^=\active
\gdef_#1{\sb{\scriptscriptstyle #1}}
\gdef^#1{\sp{\scriptscriptstyle\lycee@moinsepais #1}}
\endgroup
\newcommand\lycee@moins{\mathbin{%
  \setbox\z@\hbox{$\m@th\scriptscriptstyle\lycee@moinsorig$}%
  \dimen@=.55\wd\z@ \advance\dimen@-.5pt
  \hbox to.75\wd\z@{\hss$\m@th\scriptscriptstyle\vcenter{\hbox{%
    \tikz\draw[line width=.5pt,line cap=round](0,0)--(\the\dimen@,0);}}$\hss}}}
\begingroup\lccode`\~=`\- \lowercase{\endgroup\let~}\lycee@moins
\newcommand\lycee@moinsepais{\mathcode`\-="8000 }
\AtBeginDocument{\mathchardef\lycee@moinsorig=\mathcode`\-\relax}
\AtBeginDocument{\catcode`\_=12 \mathcode`\_="8000
  \catcode`\^=12 \mathcode`\^="8000 }
\newcommand\lycee@exposants{%
  \fontdimen13\textfont2=.48\fontdimen6\textfont2
  \fontdimen14\textfont2=.48\fontdimen6\textfont2
  \fontdimen15\textfont2=.48\fontdimen6\textfont2
  \fontdimen13\scriptfont2=.48\fontdimen6\scriptfont2
  \fontdimen14\scriptfont2=.48\fontdimen6\scriptfont2
  \fontdimen15\scriptfont2=.48\fontdimen6\scriptfont2}
\every@math@size\expandafter{\the\every@math@size\lycee@exposants}
\newlength{\retraitcalculs}
\setlength{\retraitcalculs}{2em}

\RequirePackage{tikz}
\usetikzlibrary{arrows.meta,calc,babel,shapes.misc}
\RequirePackage[most]{tcolorbox}
\tcbuselibrary{skins,breakable}

\definecolor{coulDef}{HTML}{1F4E79}
\definecolor{coulProp}{HTML}{7030A0}
\definecolor{coulRem}{HTML}{C55A11}
\definecolor{coulEx}{HTML}{2E7D32}
\definecolor{coulExo}{HTML}{B03A2E}
\definecolor{coulPart}{HTML}{1F4E79}
\definecolor{coulMeth}{HTML}{1B6E4F}
\definecolor{coulCourbe}{HTML}{0B5ED7}
\definecolor{coulCourbeB}{HTML}{C00000}
\definecolor{coulCourbeC}{HTML}{2E7D32}
\definecolor{coulGrille}{HTML}{8C8C8C}
\definecolor{coulRep}{HTML}{1B6E4F}
\definecolor{coulBar}{HTML}{2E7D32}

\newcommand{\R}{\mathbb{R}}
\newcommand{\Rs}{\mathbb{R}^{*}}
\renewcommand{\le}{\leqslant}
\renewcommand{\ge}{\geqslant}
\renewcommand{\approx}{\simeq}
\newcommand{\vect}[1]{\overrightarrow{#1}}
\newcommand{\norme}[1]{\left\lVert #1 \right\rVert}
\newcommand{\intervalle}[2]{\left[#1\,;#2\right]}
\RequirePackage{cancel}
\renewcommand{\CancelColor}{\color{coulExo}}
\newcommand{\fois}[1]{\times{\color{coulExo}#1}}
\newcommand{\barre}[1]{\cancel{#1}}

\tikzset{grille/.style={coulGrille,line width=0.4pt}}
\newcommand{\quadrillage}[4]{\draw[grille,step=1] (#1,#2) grid (#3,#4);}
\tikzset{
  croix/.style={cross out,draw,line width=0.6pt,inner sep=0pt,minimum size=4pt},
  croixdroite/.style={croix,rotate=45,line width=0.8pt},
}
\newcommand{\pt}[3]{\path (#1) node[croix]{} node[#2,font=\small,inner sep=2.5pt]{$#3$};}

\newcommand{\lycee@repere}[4]{%
  \draw[grille,step=1] (#1,#2) grid (#3,#4);
  \draw[-{Stealth[length=2mm]},thick] (#1,0)--({#3+0.4},0) node[below left=-1pt]{$x$};
  \draw[-{Stealth[length=2mm]},thick] (0,#2)--(0,{#4+0.4}) node[below left=-1pt]{$y$};
  \node[below left,font=\scriptsize,inner sep=1.5pt] at (0,0) {$0$};}
\newcommand{\lycee@gradx}[1]{\foreach \k/\l in {#1}{%
  \draw (\k,0.07)--(\k,-0.07) node[below,font=\scriptsize,inner sep=1.5pt]{$\l$};}}
\newcommand{\lycee@grady}[1]{\foreach \k/\l in {#1}{%
  \draw (0.07,\k)--(-0.07,\k) node[left,font=\scriptsize,inner sep=1.5pt]{$\l$};}}
\newcommand{\lycee@intff}[2]{\left[#1\,;#2\right]}
\newcommand{\lycee@intfo}[2]{\left[#1\,;#2\right[}
\newcommand{\lycee@intof}[2]{\left]#1\,;#2\right]}
\newcommand{\lycee@intoo}[2]{\left]#1\,;#2\right[}
\newcommand{\lycee@ens}[1]{\left\{#1\right\}}
\AtBeginDocument{%
  \providecommand{\repere}{\lycee@repere}%
  \providecommand{\gradx}{\lycee@gradx}%
  \providecommand{\grady}{\lycee@grady}%
  \providecommand{\Cf}{\mathcal{C}\sb{\scriptscriptstyle f}}%
  \providecommand{\Cg}{\mathcal{C}\sb{\scriptscriptstyle g}}%
  \providecommand{\intff}{\lycee@intff}%
  \providecommand{\intfo}{\lycee@intfo}%
  \providecommand{\intof}{\lycee@intof}%
  \providecommand{\intoo}{\lycee@intoo}%
  \providecommand{\ens}{\lycee@ens}}

\newlist{questions}{enumerate}{3}
\setlist[questions,1]{label=\arabic*\textdegree),leftmargin=*,itemsep=2pt,topsep=3pt}
\setlist[questions,2]{label=\alph*),leftmargin=*,itemsep=1pt,topsep=2pt}
\setlist[questions,3]{label=\roman*),leftmargin=*,itemsep=1pt,topsep=1pt}
\setlist[itemize]{label=\textbullet,leftmargin=*,itemsep=2pt,topsep=3pt}

\newcommand{\besoinplace}[1]{\par\penalty\@M\begingroup
  \dimen@=#1\relax\dimen@ii=\pagegoal\advance\dimen@ii-\pagetotal
  \ifdim\dimen@>\dimen@ii\ifdim\dimen@ii>\z@\vfil\fi\break\fi\endgroup}

\newcommand{\lycee@pied}{\thepage/\pageref{LastPage}}
\newcommand{\entete}[2]{%
  \pagestyle{fancy}\fancyhf{}%
  \fancyhead[L]{\footnotesize\color{coulPart}\textbf{#1}}%
  \fancyhead[R]{\footnotesize\color{coulPart}\textbf{#2}}%
  \fancyfoot[C]{\footnotesize\lycee@pied}%
  \renewcommand{\headrulewidth}{0.4pt}%
  \renewcommand{\headrule}{\hbox to\headwidth{%
    \color{coulPart!40}\leaders\hrule height \headrulewidth\hfill}}}

\RequirePackage{ccicons}
\newcommand{\lycee@licence}{{\scriptsize\color{gray}%
  \copyright\ Luc Destombe --- \ccbyncsa\ CC BY-NC-SA 4.0}}
\newif\iflycee@licence \lycee@licencetrue
\newcommand{\sanslicence}{\lycee@licencefalse}
\AtBeginDocument{\iflycee@licence\AddToHookNext{shipout/foreground}{%
  \put(0,0){\rlap{\kern\dimexpr1in+\hoffset+\oddsidemargin+\textwidth\relax
    \llap{\raisebox{-\dimexpr1in+\voffset+\topmargin+\headheight+\headsep
      +\textheight+\footskip\relax}{\lycee@licence}}}}}\fi}

\newcommand{\formulesserrees}{%
  \setlength{\abovedisplayskip}{3pt}\setlength{\belowdisplayskip}{3pt}%
  \setlength{\abovedisplayshortskip}{2pt}\setlength{\belowdisplayshortskip}{3pt}}

\setlength{\parindent}{0pt}

\RequirePackage[implicit=false,hidelinks,pdfencoding=auto,psdextra,bookmarksopen,
  bookmarksopenlevel=1,pdfcreator={},pdfproducer={}]{hyperref}
\RequirePackage{bookmark}
\hypersetup{pdfauthor={Luc Destombe},pdfkeywords={CC BY-NC-SA 4.0}}
\newcounter{lycee@signet}
\newcommand{\signet}[3][]{\stepcounter{lycee@signet}%
  \edef\lycee@dest{\if\relax\detokenize{#1}\relax
    signet.\arabic{lycee@signet}\else#1\fi}%
  \hypertarget{\lycee@dest}{}\bookmark[level=#2,dest=\lycee@dest]{#3}}
\pdfstringdefDisableCommands{%
  \def\R{R}\def\vect#1{#1}\def\Cf{Cf}\def\Cg{Cg}\def\fois#1{×#1}%
  \def\barre#1{#1}\def\intervalle#1#2{[#1;#2]}\def\intff#1#2{[#1;#2]}%
  \def\textsuperscript#1{#1}\def\dfrac#1#2{#1/#2}\def\frac#1#2{#1/#2}%
  \def\vec#1{#1⃗}}%
\RequirePackage[official]{eurosym}

\geometry{top=0.8cm,bottom=1.3cm,left=1cm,right=1cm,footskip=0.6cm}

\pagestyle{fancy}
\fancyhf{}
\renewcommand{\headrulewidth}{0pt}
\fancyfoot[C]{\footnotesize Page \thepage{} / \pageref{LastPage}}

\newcommand{\titrefiche}[2]{%
  \begin{center}
    \Large\textbf{#1}\\[2pt]
    \large\textbf{#2}\\[2pt]
  \end{center}
  \vspace{0.10cm}}

\setlength{\columnseprule}{0.4pt}
\setlength{\columnsep}{15pt}
\renewcommand{\columnseprulecolor}{\color{black!45}}
\setlength{\multicolsep}{2pt}

\newcounter{partiefiche}
\renewcommand{\thepartiefiche}{\Roman{partiefiche}}
\newif\ifapresbandeau
\newcommand{\lycee@bandeau}[1]{%
  \par\penalty-600
  \vspace{0.08cm}%
  \noindent\colorbox{coulPart}{%
    \parbox{\dimexpr\linewidth-2\fboxsep\relax}{%
      \color{white}\bfseries\raggedright\strut#1\strut}}%
  \nopagebreak\par\nopagebreak
  \vspace{0.06cm}\nopagebreak
  \apresbandeautrue}
\newcommand{\partiefiche}[1]{%
  \refstepcounter{partiefiche}%
  \lycee@bandeau{\signet{1}{\thepartiefiche{} --- #1}\thepartiefiche{} --- #1}}
\newcommand{\partiefichesansnum}[1]{\lycee@bandeau{\signet{1}{#1}#1}}

\newcounter{exo}
\newlength{\espaceexo}\setlength{\espaceexo}{7pt}
\newlength{\espacetitre}\setlength{\espacetitre}{2pt}
\newcommand{\exo}[1][\relax]{%
  \par
  \ifapresbandeau\apresbandeaufalse\else\penalty-150\addvspace{\espaceexo}\fi
  \refstepcounter{exo}%
  \noindent\signet[exercice-\arabic{exo}]{2}{Exercice \arabic{exo}}%
  \textbf{\color{coulExo}\underline{Exercice \theexo}}%
  \ifx\relax#1\relax\else\textbf{ : #1}\fi
  \par\nopagebreak\vspace{\espacetitre}\nopagebreak
  \ignorespaces}

\setlist[enumerate]{nosep,leftmargin=*,itemsep=2pt,topsep=2pt}
\setlist[itemize]{nosep,leftmargin=*,topsep=2pt}
\newcommand{\choix}[3]{\par\hspace*{1em}%
  \makebox[0.3\linewidth][l]{a) #1}\makebox[0.3\linewidth][l]{b) #2}c) #3}
\newcommand{\deux}[2]{\par\hspace*{0.6em}%
  \makebox[0.48\linewidth][l]{#1}#2}
\newcommand{\trois}[3]{\par\hspace*{0.6em}%
  \makebox[0.32\linewidth][l]{#1}\makebox[0.32\linewidth][l]{#2}#3}

  \renewcommand{\exo}[1][\relax]{%
    \ifapresbandeau\apresbandeaufalse\else\par\penalty-150\fi
    \refstepcounter{exo}%
    \vspace{0.05cm}%
    \noindent\signet[exercice-\arabic{exo}]{2}{Exercice \arabic{exo}}%
    \textbf{\color{coulExo}\underline{Exercice \theexo}}%
    \ifx\relax#1\relax\else\textbf{ : #1}\fi%
    \par\nopagebreak
    \ignorespaces}
  \newcommand{\rappel}[1]{%
    \par\vspace{0.06cm}\nopagebreak
    \noindent\fcolorbox{black!35}{black!5}{%
      \parbox{\dimexpr\linewidth-2\fboxsep-2\fboxrule\relax}{%
        \small\textbf{Rappel.} #1}}%
    \par\vspace{0.06cm}\nopagebreak}
  \setlist[enumerate]{nosep,leftmargin=*}
  \setlist[itemize]{nosep,leftmargin=*}
  \setlength{\multicolsep}{3pt plus 1pt minus 1pt}
  \setlength{\premulticols}{10pt}
  \setlength{\postmulticols}{10pt}
  \newlist{expr}{enumerate}{1}
  \setlist[expr]{label={\Alph*\ $=$},nosep,leftmargin=*,itemsep=0pt}

\renewenvironment{center}{\par\addvspace{3pt}\centering}{\par\addvspace{3pt}}
\formulesserrees
\AtBeginDocument{\formulesserrees}
\clubpenalty=10000
\widowpenalty=10000
\displaywidowpenalty=10000
\brokenpenalty=10000
\setlength{\parskip}{0pt}

\RequirePackage{ragged2e}
\setlength{\RaggedRightRightskip}{0pt plus 1fil}
\AtBeginDocument{\RaggedRight\binoppenalty=10000 \relpenalty=10000 }
\g@addto@macro\@parboxrestore{\RaggedRight}
\makeatother
\usepackage{tkz-tab}
\geometry{top=0.8cm,bottom=1.3cm,left=1cm,right=1cm,footskip=0.7cm,headheight=12pt}

\raggedcolumns
\newcommand{\qcm}[4]{\par\hspace*{0.6em}\makebox[0.48\linewidth][l]{a) #1}b) #2
  \par\hspace*{0.6em}\makebox[0.48\linewidth][l]{c) #3}d) #4}

\begin{document}

\titrefiche{1\textsuperscript{re} mathématiques spécifiques --- Calcul littéral et second
degré}{Fiche d'exercices du chapitre}

\begin{multicols}{2}

\partiefiche{Développer}

\exo[Distributivité simple]
Développer et réduire.
\trois{$A=5(x+3)$}{$B=4(2x-7)$}{$C=-3(x+6)$}
\trois{$D=-2(4-5x)$}{$E=7(3-x)$}{$F=-(x-9)$}

\exo[Un facteur en $x$]
Développer et réduire.
\trois{$A=x(x+8)$}{$B=2x(3x-1)$}{$C=-x(x-4)$}
\deux{$D=-4x(2x+5)$}{$E=3x(5-2x)$}

\exo[Double distributivité]
Développer et réduire.
\deux{$A=(x+2)(x+5)$}{$B=(x-3)(x+4)$}
\deux{$C=(x-6)(x-1)$}{$D=(2x+1)(x+3)$}
\deux{$E=(3x-2)(x+2)$}{$F=(x+5)(4-x)$}

\exo[Vrai ou faux ?]
Pour chaque égalité, dire si elle est vraie pour tout nombre $x$ ; corriger celles
qui sont fausses.
\begin{enumerate}[label=\alph*)]
  \item $3(x+4)=3x+4$
  \item $-2(x-5)=-2x+10$
  \item $x(x+1)=x^{2}+x$
  \item $(x+1)(x+2)=x^{2}+2$
  \item $-(3-x)=-3-x$
\end{enumerate}

\partiefiche{Identités remarquables}

\exo[Le carré d'une somme]
Développer et réduire.
\deux{$A=(x+4)^{2}$}{$B=(x+1)^{2}$}
\deux{$C=(2x+3)^{2}$}{$D=(3x+5)^{2}$}

\exo[Le carré d'une différence]
Développer et réduire.
\deux{$A=(x-3)^{2}$}{$B=(x-10)^{2}$}
\deux{$C=(2x-1)^{2}$}{$D=(5x-2)^{2}$}

\exo[Somme par différence]
Développer et réduire.
\deux{$A=(x+6)(x-6)$}{$B=(x-1)(x+1)$}
\deux{$C=(3x+2)(3x-2)$}{$D=(4x-5)(4x+5)$}

\exo[La bonne réponse]
Pour chaque question, une seule réponse est exacte.
\begin{enumerate}[label=\arabic*.]
  \item $(x-7)^{2}$ est égal à :
        \qcm{$x^{2}-49$}{$x^{2}-14x+49$}{$x^{2}+49$}{$x^{2}-7x+49$}
  \item $(3x+1)^{2}$ est égal à :
        \qcm{$9x^{2}+1$}{$3x^{2}+6x+1$}{$9x^{2}+6x+1$}{$9x^{2}+3x+1$}
  \item $x^{2}-16$ est le développement de :
        \qcm{$(x-4)^{2}$}{$(x-4)(x+4)$}{$(x-16)(x+1)$}{$(x-8)(x+2)$}
  \item $(2x-5)(2x+5)$ est égal à :
        \qcm{$4x^{2}-25$}{$2x^{2}-25$}{$4x^{2}-20x+25$}{$4x^{2}+25$}
\end{enumerate}

\partiefiche{Équations du premier degré}

\exo[Premières équations]
Résoudre.
\deux{a)~$x+7=12$}{b)~$4x=-20$}
\deux{c)~$3x-5=10$}{d)~$5x+2=3x+10$}
\deux{e)~$8x-1=2x+11$}{f)~$2x+9=6x-3$}

\exo[Avec des parenthèses]
Résoudre.
\deux{a)~$3(x+2)=21$}{b)~$2(x-4)=x+1$}
\deux{c)~$5(x+1)=2(x+7)$}{d)~$-(x-6)=2x$}
\deux{e)~$4(2x-1)=3(x+7)$}{}

\exo[L'inconnue au dénominateur]
Résoudre (le nombre $x$ n'est pas nul).
\[
  \text{a) }\dfrac{12}{x}=3 \qquad \text{b) }\dfrac{30}{x}=6 \qquad
  \text{c) }\dfrac{8}{x}=-2 \qquad \text{d) }\dfrac{21}{x}=7
\]
e)~Une cagnotte de $120$~€ est partagée également entre $x$ amis : chacun reçoit
$15$~€. Écrire une équation, puis trouver $x$.

\exo[Deux formules pour la piscine]
Une piscine propose l'entrée à $6$~€, ou une carte annuelle à $40$~€ avec laquelle
chaque entrée coûte $4$~€. On note $x$ le nombre d'entrées dans l'année.
\begin{enumerate}[label=\arabic*.]
  \item Calculer le prix payé pour $10$ entrées avec chaque formule.
  \item Exprimer en fonction de $x$ le prix payé avec chaque formule.
  \item Pour combien d'entrées les deux formules coûtent-elles autant ?
  \item Quelle formule choisir pour $25$ entrées ?
\end{enumerate}

\partiefiche{Équations $x^{2}=a$}

\exo[Résoudre $x^{2}=a$]
Résoudre.
\trois{a)~$x^{2}=16$}{b)~$x^{2}=1$}{c)~$x^{2}=0$}
\trois{d)~$x^{2}=3$}{e)~$x^{2}=-25$}{f)~$x^{2}=144$}

\exo[Isoler $x^{2}$ d'abord]
Résoudre.
\trois{a)~$x^{2}-9=0$}{b)~$x^{2}+4=40$}{c)~$2x^{2}=98$}
\trois{d)~$5x^{2}-20=0$}{e)~$x^{2}+10=1$}{f)~$4x^{2}-1=99$}

\exo[Des carrés]
\begin{enumerate}[label=\arabic*.]
  \item Un potager carré a une aire de $81~\text{m}^{2}$. Quelle est la longueur de
        son côté ?
  \item Une photo carrée a une aire de $100~\text{cm}^{2}$. Quelle est la longueur de
        son côté ?
  \item Quels sont les nombres dont le carré vaut $49$ ?
\end{enumerate}

\partiefiche{Équations produit}

\exo[Produit nul]
Résoudre.
\deux{a)~$(x-1)(x+7)=0$}{b)~$(x+2)(x+5)=0$}
\deux{c)~$(x-3)(x-8)=0$}{d)~$x(x+4)=0$}

\exo[Des coefficients devant $x$]
Résoudre.
\deux{a)~$(2x-4)(x+3)=0$}{b)~$(3x+9)(x-5)=0$}
\deux{c)~$(5x-15)(2x+8)=0$}{d)~$(-x+6)(4x-8)=0$}
\deux{e)~$(-2x+10)(3x+3)=0$}{}

\exo[Factoriser par $x$ d'abord]
Résoudre en factorisant par $x$.
\trois{a)~$x^{2}-4x=0$}{b)~$x^{2}+6x=0$}{c)~$3x^{2}-9x=0$}
\deux{d)~$2x^{2}+10x=0$}{e)~$x^{2}=7x$ (tout passer à gauche)}

\partiefiche{Fonction polynôme du second degré}

\exo[Reconnaître, lire $a$, $b$ et $c$]
Pour chaque fonction, dire si elle est du second degré et, si oui, donner $a$, $b$
et $c$.
\deux{$f(x)=2x^{2}+5x-1$}{$g(x)=-x^{2}+3$}
\deux{$h(x)=7x-2$}{$k(x)=4x-x^{2}$}
\deux{$m(x)=x^{2}$}{$p(x)=6-2x+3x^{2}$}

\exo[Calculer des images]
\begin{enumerate}[label=\arabic*.]
  \item Soit $f(x)=x^{2}-3x+2$. Calculer $f(0)$, $f(1)$, $f(-1)$, $f(3)$ et $f(-2)$.
  \item Soit $g(x)=-2x^{2}+x+5$. Calculer $g(0)$, $g(2)$ et $g(-1)$.
\end{enumerate}

\exo[Sur la courbe ou non ?]
Soit $f(x)=x^{2}+2x-3$. Pour chaque point $A(1\,;0)$, $B(-1\,;-2)$, $C(2\,;5)$,
$D(-3\,;0)$ et $E(0\,;3)$, dire s'il est sur la courbe de $f$.

\exo[Distance d'arrêt]
La distance d'arrêt d'une voiture, en mètres, est modélisée par $d(v)=v^{2}+3v$, où
$v$ est la vitesse en dizaines de km/h (pour $50$~km/h, $v=5$).
\begin{enumerate}[label=\arabic*.]
  \item Expliquer pourquoi $d$ est une fonction du second degré ; donner $a$, $b$ et
        $c$.
  \item Calculer $d(3)$, $d(5)$ et $d(9)$. Interpréter $d(5)$.
  \item Calculer $d(6)$. Quand la vitesse double de $30$ à $60$~km/h, la distance
        d'arrêt double-t-elle ?
\end{enumerate}

\partiefiche{Forme factorisée et racines}

\exo[Vérifier une forme factorisée]
Pour chaque fonction, vérifier l'égalité en développant la forme factorisée.
\begin{enumerate}[label=\alph*)]
  \item $f(x)=x^{2}+2x-8=(x-2)(x+4)$
  \item $g(x)=3x^{2}-3x-18=3(x+2)(x-3)$
  \item $h(x)=-x^{2}+6x-5=-(x-1)(x-5)$
  \item $k(x)=2x^{2}-8=2(x-2)(x+2)$
\end{enumerate}

\exo[Lire les racines]
Donner les racines de chaque fonction.
\deux{$f(x)=(x-3)(x+1)$}{$g(x)=2(x+5)(x-2)$}
\deux{$h(x)=-(x-4)(x-6)$}{$k(x)=5x(x-7)$}
\deux{$m(x)=-3(x+2)(x+8)$}{}

\exo[Points sur les axes]
Soit $f(x)=x^{2}-x-6$.
\begin{enumerate}[label=\arabic*.]
  \item Vérifier que $f(x)=(x+2)(x-3)$.
  \item En déduire les racines de $f$.
  \item Donner les points où la courbe de $f$ coupe l'axe des abscisses.
  \item Calculer $f(0)$. Où la courbe coupe-t-elle l'axe des ordonnées ?
\end{enumerate}

\partiefiche{Représentation graphique}

\exo[Tableau de valeurs et tracé]
Pour chaque fonction, compléter le tableau de valeurs, tracer la courbe dans un
repère, puis lire son sommet et son axe de symétrie.
\begin{enumerate}[label=\alph*)]
  \item $f(x)=x^{2}-4x+3$
\end{enumerate}
\begin{center}\small
\renewcommand{\arraystretch}{1.3}
\begin{tabular}{|c|*{7}{c|}}
\hline
$x$ & $-1$ & $0$ & $1$ & $2$ & $3$ & $4$ & $5$\\
\hline
$f(x)$ & & & & & & & \\
\hline
\end{tabular}
\end{center}
\begin{enumerate}[label=\alph*),start=2]
  \item $g(x)=-x^{2}+2x+3$
\end{enumerate}
\begin{center}\small
\renewcommand{\arraystretch}{1.3}
\begin{tabular}{|c|*{7}{c|}}
\hline
$x$ & $-2$ & $-1$ & $0$ & $1$ & $2$ & $3$ & $4$\\
\hline
$g(x)$ & & & & & & & \\
\hline
\end{tabular}
\end{center}

\exo[Associer une fonction à sa courbe]
Associer chaque fonction à sa courbe, en justifiant.
\deux{$f(x)=x^{2}-3$}{$g(x)=-x^{2}+4$}
\deux{$h(x)=2x^{2}$}{$k(x)=-x^{2}-1$}
\begin{center}
\begin{tikzpicture}[x=0.45cm,y=0.45cm]
  \repere{-4}{-5}{4}{5}
  \gradx{-3,-2,-1,1,2,3} \grady{-4,-2,2,4}
  \draw[coulCourbe,line width=1.1pt,domain=-3:3,samples=60] plot (\x,{-\x*\x+4});
  \draw[coulCourbeB,line width=1.1pt,domain=-1.58:1.58,samples=60] plot (\x,{2*\x*\x});
  \draw[coulCourbeC,line width=1.1pt,domain=-2.8:2.8,samples=60] plot (\x,{\x*\x-3});
  \draw[black!70,line width=1.1pt,domain=-2:2,samples=60] plot (\x,{-\x*\x-1});
  \node[coulCourbe,font=\footnotesize,fill=white,inner sep=1pt] at (-3.4,-4.3) {$\mathcal{C}_{1}$};
  \node[coulCourbeB,font=\footnotesize,fill=white,inner sep=1pt] at (2.1,4.6) {$\mathcal{C}_{2}$};
  \node[coulCourbeC,font=\footnotesize,fill=white,inner sep=1pt] at (3.4,4.3) {$\mathcal{C}_{3}$};
  \node[black!70,font=\footnotesize,fill=white,inner sep=1pt] at (-2.6,-3.3) {$\mathcal{C}_{4}$};
\end{tikzpicture}
\end{center}

\exo[Lire une parabole]
La courbe de $f(x)=x^{2}+2x-3$ est tracée ci-dessous.
\begin{center}
\begin{tikzpicture}[x=0.45cm,y=0.45cm]
  \repere{-5}{-5}{3}{6}
  \gradx{-4,-3,-2,-1,1,2} \grady{-4,-3,-2,-1,1,2,3,4,5}
  \draw[coulCourbe,line width=1.1pt,domain=-4:2,samples=60] plot (\x,{\x*\x+2*\x-3});
  \node[coulCourbe,font=\footnotesize] at (2.4,4) {$\Cf$};
\end{tikzpicture}
\end{center}
\begin{enumerate}[label=\arabic*.]
  \item La parabole est-elle tournée vers le haut ou vers le bas ? Est-ce cohérent
        avec l'expression de $f$ ?
  \item Lire les coordonnées du sommet et l'axe de symétrie.
  \item Lire les racines de $f$.
  \item Résoudre graphiquement $f(x)=-3$.
  \item Calculer $f(0)$ et retrouver ce point sur la courbe.
\end{enumerate}

\partiefiche{Tableaux de variations}

\exo[Sommet et tableau de variations]
Pour chaque fonction, donner le sommet de la parabole, puis le tableau de
variations.
\deux{a)~$f(x)=(x-2)(x-6)$}{b)~$g(x)=-2(x+1)(x-3)$}
\deux{c)~$h(x)=3(x+4)(x+2)$}{d)~$k(x)=-(x-5)(x+1)$}

\exo[Factoriser d'abord]
Pour chaque fonction, trouver les racines, puis le sommet et le tableau de
variations.
\deux{a)~$f(x)=x^{2}-6x$}{b)~$g(x)=-x^{2}+10x$}
\deux{c)~$h(x)=2x^{2}+8x$}{d)~$k(x)=x^{2}-9$}

\exo[Maximum ou minimum ?]
\begin{enumerate}[label=\arabic*.]
  \item Pour chaque fonction $f$, $g$ et $h$, dire si elle admet un maximum ou un
        minimum, et le donner :
        \par $f(x)=4(x-1)(x-3)$ \quad $g(x)=-(x+3)(x-1)$ \quad $h(x)=-5x(x-2)$.
  \item Calculer $f(0)$ et $f(4)$, puis dresser le tableau de variations de $f$
        entre $0$ et $4$.
\end{enumerate}

\partiefiche{Problèmes type épreuve anticipée}

\exo[Un food truck]
Un food truck vend entre $0$ et $120$ repas par jour. On note $x$ le nombre de
repas, en dizaines ($x$ entre $0$ et $12$). Le bénéfice de la journée, en dizaines
d'euros, est $B(x)=-x^{2}+12x-20$.
\begin{enumerate}[label=\arabic*.]
  \item Calculer $B(0)$. Interpréter ce résultat dans le contexte de l'exercice.
  \item Vérifier que $B(x)=-(x-2)(x-10)$.
  \item Pour combien de repas le bénéfice est-il nul ?
  \item Dresser le tableau de variations de $B$ entre $0$ et $12$.
  \item Quel est le bénéfice maximal ? Pour combien de repas ?
\end{enumerate}

\exo[Une balle lancée d'une falaise]
On lance une balle du haut d'une falaise. Sa hauteur au-dessus de la mer, en
mètres, au bout de $t$ secondes est $h(t)=-5t^{2}+10t+15$.
\begin{enumerate}[label=\arabic*.]
  \item Calculer $h(0)$. Interpréter.
  \item Vérifier que $h(t)=-5(t+1)(t-3)$.
  \item Au bout de combien de temps la balle touche-t-elle l'eau ?
  \item Dresser le tableau de variations de $h$ entre $0$ et $3$.
  \item Quelle hauteur maximale la balle atteint-elle ? À quel instant ?
\end{enumerate}

\exo[Le plus grand enclos]
Avec $40$~m de grillage, on entoure un enclos rectangulaire. On note $x$ la longueur
d'un côté, en mètres ($x$ entre $0$ et $20$).
\begin{enumerate}[label=\arabic*.]
  \item Si $x=5$, quelle est la longueur de l'autre côté ? Quelle est l'aire ?
  \item Justifier que l'autre côté mesure $20-x$ et que l'aire est
        $A(x)=x(20-x)$.
  \item Développer $A(x)$, puis donner les racines de $A$.
  \item Dresser le tableau de variations de $A$ entre $0$ et $20$.
  \item Quelles dimensions donnent la plus grande aire ? Quelle est-elle ?
\end{enumerate}

\partiefichesansnum{Exercices type épreuve anticipée}

\exo[Automatismes (QCM)]
Pour chaque question, une seule réponse est exacte.
\begin{enumerate}[label=\arabic*.]
  \item La solution de l'équation $6x+1=2x+9$ est :
        \qcm{$1$}{$2$}{$4$}{$10$}
  \item $(x+3)^{2}$ est égal à :
        \qcm{$x^{2}+9$}{$x^{2}+6x+9$}{$x^{2}+3x+9$}{$x^{2}+6x+6$}
  \item L'ensemble des solutions de $x^{2}=7$ est :
        \qcm{$\ens{7}$}{$\ens{-\sqrt{7}\,;\sqrt{7}}$}{$\ens{\sqrt{7}}$}{$\ens{-3{,}5\,;3{,}5}$}
  \item L'ensemble des solutions de $(3x-6)(x+4)=0$ est :
        \qcm{$\ens{-2\,;4}$}{$\ens{-4\,;2}$}{$\ens{-4\,;3}$}{$\ens{-4\,;6}$}
  \item Le point d'abscisse $2$ de la courbe d'équation $y=x^{2}-3x+1$ est :
        \qcm{$(2\,;-1)$}{$(2\,;3)$}{$(2\,;11)$}{$(2\,;-3)$}
  \item La fonction $f(x)=-2(x-1)(x-5)$ admet :
        \qcm{un minimum en $3$}{un maximum en $3$}{un maximum en $1$}{un minimum en $5$}
\end{enumerate}

\exo[Une épidémie de grippe]
Dans une ville, le nombre de malades de la grippe, en centaines, $x$ semaines après
le début de l'épidémie ($x$ entre $0$ et $8$), est modélisé par $f(x)=-x^{2}+8x$.
\begin{center}
\begin{tikzpicture}[x=0.6cm,y=0.2cm]
  \draw[grille,xstep=1,ystep=2] (0,0) grid (9,18);
  \draw[-{Stealth[length=2mm]},thick] (0,0)--(9.4,0) node[below left=-1pt]{$x$};
  \draw[-{Stealth[length=2mm]},thick] (0,0)--(0,19.2) node[below left=-1pt]{$y$};
  \node[below left,font=\scriptsize,inner sep=1.5pt] at (0,0) {$0$};
  \foreach \i in {1,...,8} {\draw (\i,0.3)--(\i,-0.3) node[below,font=\scriptsize,inner sep=1.5pt]{$\i$};}
  \foreach \v in {2,4,...,16} {\draw (0.07,\v)--(-0.07,\v) node[left,font=\scriptsize,inner sep=1.5pt]{$\v$};}
  \draw[coulCourbe,line width=1.1pt,domain=0:8,samples=60] plot (\x,{-\x*\x+8*\x});
  \node[coulCourbe,font=\footnotesize] at (7.6,11) {$\Cf$};
\end{tikzpicture}
\end{center}
\begin{enumerate}[label=\arabic*.]
  \item Lire $f(2)$ sur le graphique et interpréter.
  \item Déterminer graphiquement la semaine du pic de l'épidémie et le nombre de
        malades cette semaine-là.
  \item Calculer $f(1)$. Interpréter ce résultat dans le contexte de l'exercice.
  \item Vérifier que $f(x)=-x(x-8)$. Au bout de combien de semaines l'épidémie
        s'arrête-t-elle ?
  \item Dresser le tableau de variations de $f$ entre $0$ et $8$ et retrouver par le
        calcul le nombre de malades au pic.
\end{enumerate}

\exo[Le prix d'une place de théâtre]
Une troupe de théâtre fixe le prix d'une place à $x$ euros ($x$ entre $0$ et $20$).
Elle prévoit alors $200-10x$ spectateurs.
\begin{enumerate}[label=\arabic*.]
  \item Pour une place à $8$~€, combien de spectateurs la troupe prévoit-elle ? Quelle
        est la recette ?
  \item Quel prix donne $150$ spectateurs ?
  \item Justifier que la recette est $R(x)=x(200-10x)$, puis que
        $R(x)=-10x^{2}+200x$.
  \item Donner les racines de $R$ et les interpréter.
  \item Dresser le tableau de variations de $R$ entre $0$ et $20$.
  \item Quel prix donne la plus grande recette ? Quelle est-elle ?
\end{enumerate}

\end{multicols}
\end{document}
